Los controles son los del arcade. El teclado, los cursores y el joystick 1 funcionan a la vez.

\begin{center}
\begin{tabular}{@{}lll@{}}
\toprule
\textbf{\tAction} & \textbf{\tKeyboard} & \textbf{\tJoystick}\\
\midrule
Mover izquierda / derecha & \key{$\leftarrow$} \key{$\rightarrow$} & palanca izquierda / derecha\\
Arriba: saltar al piso de arriba & \key{$\uparrow$} & palanca arriba\\
Abajo: caer al piso de abajo & \key{$\downarrow$} & palanca abajo\\
Patada & \key{ESPACIO} o \key{Z} & bot\'on A\\
Salto & \key{M} o \key{X} & bot\'on B\\
Empezar, modo dif\'icil (t\'itulo) & \key{1}, \key{ESPACIO} & bot\'on A\\
Empezar, modo medio (t\'itulo) & \key{2} & bot\'on B\\
Letras de tus iniciales & \key{$\leftarrow$} \key{$\rightarrow$}, patada para aceptar & palanca, bot\'on A\\
\bottomrule
\end{tabular}
\end{center}

\textbf{Notas.}
\begin{itemize}
\item La palanca se comporta como la del arcade. Si cambias de sentido sin pasar por el centro, el juego le da a la palanca un fotograma neutro para que el h\'eroe se d\'e la vuelta en lugar de andar hacia atr\'as.
\item Patea mientras saltas para dar una \emph{patada voladora}, que da m\'as puntos.
\item \key{ESPACIO}, el bot\'on A y el bot\'on B empiezan la partida solo desde el t\'itulo y la demo. Durante el juego, y al escribir las iniciales, solo patean y saltan.
\item Los dos modos son partidas de un jugador. \textbf{Dif\'icil} es el ajuste duro del arcade; \textbf{medio} da a los enemigos menos velocidad y agresividad, y tambi\'en se vuelve m\'as dif\'icil con cada fase.
\item No hay tecla de pausa ni de reinicio, como en el arcade.
\end{itemize}
